\section{What To Do In Your Agent}
\label{sec:implications}

This section is the practical summary. It assumes nothing from
Section~\ref{sec:model} beyond the phrase ``reservation price''.

\subsection{A checklist}

Before adding a stopping optimizer, answer three questions about the situation your
agent is actually in.

\begin{enumerate}
  \item \textbf{Do merchants charge different prices for this product?} Measure the
    ratio of the dearest to the cheapest offer on a sample of real products. Below
    roughly \ResultNoAdvantageDispersionAtOrBelow{}$\times$, stop here and use a
    constant rule; there is nothing to optimize and the machinery will not pay for
    itself.
  \item \textbf{Can you afford to query most of the alternatives?} If your budget
    covers only a small fraction of the reachable merchants, an adaptive rule can
    make things worse by declining acceptable offers it cannot follow up. Either
    raise the budget or accept early deliberately.
  \item \textbf{Is more than one alternative genuinely reachable?} With a single
    fallback, an optimal rule and a one-step lookahead coincide. Implement the simple
    one.
\end{enumerate}

If all three point the same way, the rule is worth implementing, and
Table~\ref{tab:decision} gives the improvement to expect.

\subsection{Design guidance}

\paragraph{Show the stopping rule the budget.}
The most common failure we see is a stopping policy that knows the prices but not the
remaining resources. Remaining budget belongs in the state that decides whether to
continue. A merchant that cannot be afforded should disappear from the decision, not
be scored and then rejected.

\paragraph{Reserve before you dispatch, and charge pessimistically.}
Deduct the full permit before the call, reclaim what is provably unused afterwards,
and charge the full permit whenever the call is cancelled or truncated and the true
usage is unknown. This makes overruns impossible rather than rare, which matters more
than the small efficiency it costs.

\paragraph{Keep forecasting separate from control.}
The reservation price is a function of the forecast; it is not itself a forecast.
Keeping them apart means an improved price or availability model can be swapped in
without touching the stopping logic, and lets you attribute a regression to one or
the other.

\paragraph{Price failure explicitly.}
An episode ending without a purchase is not a zero-cost outcome, and a policy scored
as though it were will learn to stall. Give failure a number and use the same number
everywhere.

\paragraph{Measure your own dispersion before trusting ours.}
The decision table is indexed by quantities you can measure in a day. Our corpus is
one protocol at two dates and its dispersion is unusually low; yours may differ, and
if it does, the table and not our headline number is the thing to apply.
